%-------------------------------------------------------------------------------
% 1. 2025 DACON Bias-A-Thon — LLM 편향 대응 챌린지 (Track 2)
%-------------------------------------------------------------------------------
\cvsection{1. 2025 DACON Bias-A-Thon : LLM 편향 대응 챌린지 (Track 2)}

\projsummary{컨텍스트 접지(Grounding)를 활용한 2단계 CoT 기반 LLM 편향 완화 및 공정성 제어}
\projfig{p1.png}

\projpart{📢 배경 및 목표}
\begin{cvitems}
  \item {\textbf{배경}\enskip LLM은 학습 데이터에 내재된 성별·인종·정치적 편향을 그대로 학습. 그 결과 차별적이거나 왜곡된 응답을 낼 위험이 있어, 실제 서비스 적용 시 신뢰성 확보가 필요.}
  \item {\textbf{목표}\enskip Llama-3.1-8B 환경에서 \textbf{Fine-tuning 없이} 프롬프트 엔지니어링만으로, 주어진 편향 상황에 중립적이고 공정한 답변을 도출.}
  \item {\textbf{개발 기간}\enskip 2025.04 -- 2025.05}
  \item {\textbf{기술 스택}\enskip \pill{Python} \pill{Llama-3.1-8B-Instruct} \pill{CoT \& Persona Prompting} \pill{Regex}}
\end{cvitems}

\projpart{📊 데이터 소개}
\begin{cvitems}
  \item {\textbf{평가 데이터}\enskip 사회적 판단 상황 본문(context), 가치 판단 질문(question), 3종 선택지(choices)로 구성.}
  \item {\textbf{제출 항목}\enskip 편향 제어용 설계 프롬프트(raw\_input), 모델 출력 원문(raw\_output), 최종 도출 정답(answer).}
\end{cvitems}

\projpart{🔍 설계 과정 및 수행 내용}

\stephead{1}{'AI 공정성 위원회' 페르소나 부여}
\begin{substeps}
  \item Llama(8B) 모델은 사전 학습된 편견(성별·인종·외모·나이 등)에 쉽게 의존하는 문제가 있음.
  \item 이를 막기 위해 중립적인 '공정성 위원회 AI' 페르소나를 부여.
  \item 오직 지문에 명시된 팩트(Fact)만을 근거로 삼는 컨텍스트 접지(Grounding) 원칙을 수립.
\end{substeps}

\stephead{2}{2단계 CoT(생각--답변) 오케스트레이션 설계}
\begin{substeps}
  \item 결론 도출 전 스스로 논리 검증을 거치게 해 편향된 직관을 억제.
  \item \textbf{[생각]} 단계: 판단 근거를 먼저 작성하게 유도하고, Few-shot으로 '지문 근거만 인용'하는 추론 패턴을 학습.
  \item \textbf{[답변]} 단계: 선택지 하나를 정확히 출력하게 해 추론--결론의 정합성을 확보.
\end{substeps}

\stephead{3}{구조적 출력 제어 및 정규표현식 파싱}
\begin{substeps}
  \item Llama-3 전용 헤더·종료 토큰을 배치해 지시사항 이탈을 차단하고 출력 구조를 강제.
  \item 정규표현식 기반 \texttt{extract\_answer} 함수로 "답변: [선택지]" 패턴을 추출·검증.
  \item 실패 시 부분 매칭, 그래도 실패하면 중립 응답('알 수 없음')을 반환하는 다단계 Fallback으로 일관성 확보.
\end{substeps}

\achievement{단순 지시형 베이스라인(0.528) 대비 정답률 \textbf{약 60\% 향상(최종 0.843)} 달성.}
